\begin{table}[htbp]
\centering
\caption{Minimal winning coalitions by ideological domain}
\label{tab:interval-summary}
\small
\setlength{\tabcolsep}{4pt}
\begin{tabularx}{\textwidth}{@{}rcrr>{\raggedright\arraybackslash}X@{}}
\toprule
Election & \(k\) & Minimal majorities & Inversions & Strongest minimal inversion \\
\midrule
2014 & 0 & 8 & 4 & PTB--PR (47.45\%; 257 seats) \\
2014 & 1 & 87 & 46 & PTB--PR (47.45\%; 257 seats) \\
2018 & 0 & 8 & 1 & PT--PSDB (49.95\%; 267 seats) \\
2018 & 1 & 108 & 42 & PCdoB--PODE, omitting PSDB (47.59\%; 257 seats) \\
2022 & 0 & 4 & 2 & PP--PL (45.26\%; 258 seats) \\
2022 & 1 & 39 & 12 & PP--PL (45.26\%; 258 seats) \\
\bottomrule
\end{tabularx}
\begin{flushleft}
\footnotesize Notes: The ideological order contains seat-winning parties.
\(k\) counts omitted parties inside a coalition's ideological span.
Minimality is defined within each stated domain. Vote shares use all valid federal-deputy votes.
The strongest minimal inversion has the lowest national vote share; exact ties use coalition size and then canonical membership order.
\end{flushleft}
\end{table}
